\begin{thema}{A}
\begin{erwthma}
    \item \bmath{Δίνεται συναρτήση $ g $ ορισμένη σε ένα διάστημα $ [a,\beta] $ τέτοια ώστε $ g(x)\leq 0 $ για κάθε $ x\in[a,\beta] $. Να αποδείξετε ότι το εμβαδόν του χωρίου $ \varOmega $ μεταξύ της $ C_g $, του άξονα $ x'x $ και των ευθειών $ x=a,x=\beta $ δίνεται από τον τύπο \[ E(\varOmega)=-\int_{a}^{\beta}{g(x)) \d x} \]}
    \item Τι ονομάζουμε αρχική συνάρτηση ή παράγουσα μιας συνάρτησης $f$ σε ένα διάστημα $\Delta$;
    \item Πώς ορίζονται οι πράξεις της πρόσθεσης $f+g$ και του πολλαπλασιασμού $f \cdot g$ δύο συναρτήσεων; Ποιο το πεδίο ορισμού τους;
    \item Να χαρακτηρίσετε τις προτάσεις που ακολουθούν, γράφοντας στο τετράδιό σας, δίπλα στο γράμμα που αντιστοιχεί σε κάθε πρόταση, τη λέξη \textbf{Σωστό}, αν η πρόταση είναι σωστή, ή \textbf{Λάθος}, αν η πρόταση είναι λανθασμένη.
    \begin{alist}
        \item Αν $C_f$ έχει οριζόντια ασύμπτωτη στο $+\infty$, δεν μπορεί να έχει και πλάγια ασύμπτωτη στο $+\infty$.
        \item Αν $\dint_a^\beta f(x)dx > 0$ με $a < \beta$, τότε αναγκαστικά $f(x) > 0$ για κάθε $x \in [a, \beta]$.
        \item Το εμβαδόν μεταξύ δύο ευθειών που τέμνονται και του άξονα ολοκληρώνεται χωρίς να χωριστεί το διάστημα στην προβολή της τομής τους.
        \item Αν $f(x) \le 0$ στο $[a, \beta]$, τότε το εμβαδόν του χωρίου δίνεται από $-\dint_a^\beta f(x) dx$.
        \item Ο συμβολισμός της μεταβλητής ολοκλήρωσης $x$ ή $t$ δεν προκαλεί καμία απολύτως αλλαγή στην αριθμητική τιμή ενός ορισμένου ολοκληρώματος.
    \end{alist}

    % ----------------- GENERAL THEORY ----------------- %
\end{erwthma}
\end{thema}
